\begin{example}[Affine line over a field]\label{ex:vi26-r01}
The structure morphism $\mathbb A^1_k=\Spec k[t]\to\Spec k$ is of finite type because $k[t]$ is generated by $t$ as a $k$-algebra.
\end{example}

\begin{example}[Affine $n$-space]\label{ex:vi26-r02}
$\mathbb A^n_A=\Spec A[x_1,\dots,x_n]\to\Spec A$ is of finite type; the coordinate algebra has the displayed finite generating set.
\end{example}

\begin{example}[A hypersurface family]\label{ex:vi26-r03}
For $B=A[x_1,\dots,x_n]/(F)$, the morphism $\Spec B\to\Spec A$ is of finite type.  The number of equations is irrelevant to finite type; only the number of algebra generators matters.
\end{example}

\begin{example}[An arbitrary affine quotient]\label{ex:vi26-r04}
For any ideal $I\subseteq A[x_1,\dots,x_n]$, even one not finitely generated, $A[x_1,\dots,x_n]/I$ is a finite-type $A$-algebra.
\end{example}

\begin{example}[A principal localization]\label{ex:vi26-r05}
$A_f\cong A[t]/(ft-1)$ is finitely generated over $A$, so $D(f)\to\Spec A$ is of finite type.
\end{example}

\begin{example}[The multiplicative group as an affine scheme]\label{ex:vi26-r06}
$\Spec k[t,t^{-1}]\to\Spec k$ is of finite type because $k[t,t^{-1}]\cong k[t,u]/(tu-1)$.
\end{example}

\begin{example}[Polynomial algebra in infinitely many variables]\label{ex:vi26-r07}
$\Spec k[x_1,x_2,\ldots]\to\Spec k$ is not locally of finite type; infinitely many algebra generators remain near the generic point.
\end{example}

\begin{example}[Infinite disjoint union of points]\label{ex:vi26-r08}
$\coprod_{n\ge1}\Spec k\to\Spec k$ is locally of finite type but not finite type because the source is not quasi-compact.
\end{example}

\begin{example}[A finite disjoint union of points]\label{ex:vi26-r09}
$\coprod_{i=1}^r\Spec k\to\Spec k$ is finite type: it is locally finite type and the source has a finite affine cover.
\end{example}

\begin{example}[Coordinate projection]\label{ex:vi26-r10}
$\Spec k[x,y]\to\Spec k[x]$ is of finite type because $k[x,y]=k[x][y]$.
\end{example}

\begin{example}[Restriction to a basic open]\label{ex:vi26-r11}
Restricting the preceding projection over $D(x)$ gives $\Spec k[x,x^{-1},y]\to\Spec k[x,x^{-1}]$, again of finite type.
\end{example}

\begin{example}[Scalar extension]\label{ex:vi26-r12}
If $B$ is finite type over $A$ and $A\to A'$ is any ring map, then $B\otimes_AA'$ is finite type over $A'$.
\end{example}

\begin{example}[A fiber of affine space]\label{ex:vi26-r13}
The fiber of $\mathbb A^n_S\to S$ at $s$ is $\mathbb A^n_{\kappa(s)}$, a finite-type scheme over the residue field.
\end{example}

\begin{example}[The integers are Noetherian]\label{ex:vi26-r14}
$\Spec\mathbb Z$ is Noetherian because $\mathbb Z$ is a principal ideal domain, hence Noetherian.
\end{example}

\begin{example}[Affine space over the integers]\label{ex:vi26-r15}
$\mathbb Z[x_1,\dots,x_n]$ is Noetherian by Hilbert's basis theorem, so $\mathbb A^n_{\mathbb Z}$ is Noetherian.
\end{example}

\begin{example}[A distinguished open in a Noetherian affine scheme]\label{ex:vi26-r16}
If $A$ is Noetherian, then $A_f$ is Noetherian; hence every distinguished open $D(f)\subseteq\Spec A$ is Noetherian.
\end{example}

\begin{example}[A non-Noetherian affine scheme]\label{ex:vi26-r17}
$\Spec k[x_1,x_2,\ldots]$ is not Noetherian because $(x_1)\subsetneq(x_1,x_2)\subsetneq\cdots$ is a strictly ascending chain.
\end{example}

\begin{example}[Composition]\label{ex:vi26-r18}
$\Spec k[x,y]\to\Spec k[x]\to\Spec k$ is a composite of finite-type morphisms, hence finite type.
\end{example}

\begin{example}[Field extension]\label{ex:vi26-r19}
For any extension $K/k$, base changing a finite-type $k$-scheme $X$ gives the finite-type $K$-scheme $X_K$.
\end{example}

\begin{example}[Finite type over an affine Noetherian base]\label{ex:vi26-r20}
If $A$ is Noetherian and $B=A[b_1,\dots,b_n]$, then $B$ is Noetherian; thus $\Spec B\to\Spec A$ preserves local Noetherianity.
\end{example}